\section{Optical operators and their physical realization}
\label{sec:want-optics}

\begin{panel}
\textbf{What the additional source supplies.}
W1, physical pages 11--16, supplies a mathematical optical interaction chain:
aperture, propagation, spectral acquisition, material response, an event
instrument and an energy ledger. These are usable additions to the exact seed
as typed models and observations. They specify how a photonic instrument could
interact with the kernel. A photonic implementation of the kernel's arithmetic,
memory and delayed feedback remains a separate realization problem.
\end{panel}

This chapter preserves the source equations and states the conditions under
which they can be lifted into XOP-R1. Newly derived consequences are identified
as propositions. A mathematical identity within a chosen optical model is not
an experimental certificate for an unspecified material or assembled device.
No embedded program, numerical experiment or physical test was run for R12.

\subsection{A slit is an exact spatial restriction}

In the aperture plane, let $\mathbf a,\mathbf b$ be orthogonal unit directions,
$\mathbf r_c$ its centre, and $w_s,h_s>0$ its width and height. The W1 aperture is
\begin{equation}
 A(\mathbf r)=
 \mathbf1_{\{|\mathbf a\cdot(\mathbf r-\mathbf r_c)|\le w_s/2\}}
 \mathbf1_{\{|\mathbf b\cdot(\mathbf r-\mathbf r_c)|\le h_s/2\}}.
 \label{eq:want-aperture}
\end{equation}
The complete geometric expression can remain in a seed. A pixel grid is not
needed to \emph{define} its edges, centre, deformation or membership predicate.
At an exact boundary the declared $\le$ convention includes the point; uncertain
physical edge locations remain uncertain parameters. A shift changes the
centre, a tilt changes the plane frame, and strain changes the edge model.
Those are distinct state changes, even when driven by one actuator.

For a square-integrable scalar field, $(P_AU)(\mathbf r)=A(\mathbf r)U(\mathbf r)$
satisfies
\begin{equation}
 P_A^2=P_A,\qquad
 \|P_AU\|_2^2=\int A|U|^2\,dA\le\int |U|^2\,dA.
 \label{eq:want-slit-contraction}
\end{equation}
The proof uses $A^2=A$ and $0\le A\le1$ pointwise. This is the exact Boolean
projection supplied by the ideal aperture. Multiplication by a general complex
transmission $t_\lambda$ instead gives $(M_t)^2U=t_\lambda^2U$; it is idempotent
only when the relevant transmission values satisfy $t_\lambda^2=t_\lambda$.
Consequently a real absorber or phase coating is not automatically the same
operator as an ideal ASA/NA word mask.

\subsection{A declared field transfer, including diffraction}

The source interaction chain is
\begin{equation}
 U_{\mathrm{det},\lambda}
 =H_{z,\lambda}\!\left[
 A\,t_\lambda(T,a,\zeta)
 e^{i\phi_{\mathrm{lens},\lambda}}U_{\mathrm{in},\lambda}\right].
 \label{eq:want-field-transfer}
\end{equation}
Here $T$ is temperature, $a$ a material hysteresis state, and $\zeta$ a reporter
fraction. The propagation medium, polarization approximation, input/output
planes, field normalization and boundary conditions belong to $H$. One may
choose an analytically specified propagator or a calibrated finite-mode
transfer. A name, lookup-table address or trained transfer without its
coefficients and domain does not complete the physical definition.

For example, set $\lambda_m=\lambda/n_m$, $k_m=2\pi/\lambda_m$, and use local
coordinates $(x,y)$ in a centred rectangular aperture. For a bounded aperture
field $U_A$, a Fraunhofer angular shape transform is
\begin{equation}
 \mathcal F(s_x,s_y)=
 \int_{-w_s/2}^{w_s/2}\!\int_{-h_s/2}^{h_s/2}
 U_A(x,y)e^{-ik_m(s_xx+s_yy)}\,dy\,dx.
 \label{eq:want-fraunhofer}
\end{equation}
For this nested Riemann expression, each inner $y$ section of the full
integrand must be Riemann-integrable at every admitted $x$, and the resulting
inner integral must be Riemann-integrable as a function of $x$. Joint
two-dimensional integrability alone does not establish all section domains.
A continuous aperture field on the closed rectangle is one sufficient
condition; other fields need their stated section and outer-integral
conditions. The variables $s_x,s_y$ are direction cosines in the relevant far-field
directions. This expression specifies the angular shape; the physical
propagation prefactor and flux measure must still match the selected $H$.
For uniform complex $U_A=U_0$,
\begin{equation}
 \mathcal F=U_0w_sh_s\,\operatorname{sinc}(k_mw_ss_x/2)
                  \operatorname{sinc}(k_mh_ss_y/2),\qquad
 \operatorname{sinc}u=
 \begin{cases}\sin u/u,&u\ne0,\\1,&u=0.\end{cases}
 \label{eq:want-sinc}
\end{equation}
The value at zero is a declared continuous extension, not an eager division by
zero. The first-zero relations are
\begin{equation}
 \sin\alpha_x=\frac{\lambda}{n_mw_s},\qquad
 \sin\alpha_y=\frac{\lambda}{n_mh_s},
 \label{eq:want-first-zero}
\end{equation}
when the corresponding ratios are at most one and the adopted scalar
Fraunhofer assumptions apply. Narrowing a dimension broadens its diffraction
pattern. The MIT rectangular-aperture optics demonstration provides a direct
reference for this Fourier-optics dependence.
\href{https://ocw.mit.edu/courses/res-6-006-video-demonstrations-in-lasers-and-optics-spring-2008/e06ca88aa4b673d48884b8705496a1b7_4YPxRTFxy2A.pdf}{[P1]}

The integral and its factorized form can both be stored as exact terms. That
removes compulsory spatial rasterization from the definition; it does not
remove diffraction, the far-field approximation, or uncertainty in a measured
incident field. A finite modal expansion or sampled transfer is a declared
model reduction with its own omitted-field error. Physical light outside a
geometric occupancy set may still affect the detector through diffraction and
scattering.

\subsection{Passive transfer is a power condition}

With power-normalized input and output modes, a passive linear transfer $S$
must obey
\begin{equation}
 \|S\mathbf u\|_2^2\le\|\mathbf u\|_2^2
 \quad\hbox{for every }\mathbf u,
 \qquad S^\dagger S\preceq I.
 \label{eq:want-passive}
\end{equation}
More generally, positive flux metrics give
$S^\dagger W_{\mathrm{out}}S\preceq W_{\mathrm{in}}$.
An electric-field transmission coefficient across different refractive media
cannot be judged by its magnitude alone; its flux normalization matters.
For equal power conventions, the composition of the aperture, a multiplier
with $|t_\lambda|\le1$, a real phase factor, and a nonexpansive $H$ is itself
nonexpansive. Each factor bounds the next norm, which proves the claim.
Passive concentration can increase local irradiance while preserving the total
energy budget; it does not increase a fixed-frequency photon's energy.

An active stage may supply gain. Its pump, electrical power, added fluctuations,
bandwidth and saturation then belong to the physical profile. Passivity is not
an extra restriction on exact mathematical arithmetic; it is a necessary
condition only for the particular proposed passive field implementation.

\subsection{Liquid lensing and a shared spectral budget}

The W1 lens proposal becomes physical only after its interface and actuator are
specified. For a static fluid interface and a thin paraxial lens,
\begin{align}
 n_1\sin\theta_1&=n_2\sin\theta_2,\\
 \Delta p&=\gamma(1/R_1+1/R_2),\\
 \frac1{f_\lambda}
 &=\left(\frac{n_\ell(\lambda)}{n_m(\lambda)}-1\right)
   \left(\frac1{R_{\mathrm{front}}}-\frac1{R_{\mathrm{back}}}\right).
 \label{eq:want-lens}
\end{align}
The radii have one declared optical sign convention and nonzero finite values
where reciprocals are used. A plane face is represented by zero curvature,
not by a finite real literal called infinity. The refractive indices are
positive, and propagating Snell branches require their stated domains.
Contact-angle and support
conditions close the shape problem; pressure, voltage or another declared
actuator supplies its control. Berge and Peseux reported a liquid lens whose
contact angle was controlled by electrocapillarity. That is a physical
precedent for variable liquid lensing, not a calibration of the proposed rim
or matrix. The publisher abstract was inspected.
\href{https://link.springer.com/article/10.1007/s101890070029}{[P2]}

Let $P_\lambda(t)\ge0$ denote incident spectral power per positive vacuum
wavelength at a specified plane, over $t_{n+1}\ge t_n$. For nonoverlapping
passive output branches,
\begin{equation}
 0\le\tau_k(\lambda),\quad
 \tau_1(\lambda)+\tau_2(\lambda)\le1,\quad
 0\le\eta_k(\lambda)\le1,
 \label{eq:want-spectral-budget}
\end{equation}
under a photon-detection efficiency convention with at most one primary count
per delivered photon. Internal gain and acquisition codes require their own
response definition. Before dead time and saturation, the ideal mean primary
detections under that efficiency model are
\begin{equation}
 \mu_k=\int_{t_n}^{t_{n+1}}\!\int
 \eta_k(\lambda)\tau_k(\lambda)P_\lambda(t)
 \frac{\lambda}{hc}\,d\lambda\,dt,
 \quad k=1,2.
 \label{eq:want-photon-mean}
\end{equation}
Background, dead time, saturation and readout noise complete the observation
model and may change its mean as well as its distribution; $\mu_k$ is not the
particular count that will be observed. Both ideal means derive
from one incident budget. An ideally routed single photon cannot be counted
in both destructive, nonoverlapping branches. A one-bit event may use a
specified primary band or a suitable joint score without requiring both
branches to fire.

With all outgoing reflected and transmitted channels included,
\begin{equation}
 \mathcal R_\lambda+\mathcal T_\lambda+\mathcal A_\lambda=1,
 \qquad P_{\mathrm{abs}}=\int\mathcal A_\lambda P_\lambda\,d\lambda.
 \label{eq:want-absorption}
\end{equation}
These fractions describe the selected passive interaction. After absorption,
energy may reside in chemical excitation, heat or another store and may later
leave as radiation or electrical work. Multiple readouts of that event share
its energy and any supplied amplification energy.

\subsection{The material state closes the optical feedback}

The reporter and thermal derivations in Section~\ref{sec:want-energy} close
several otherwise missing arguments of $t_\lambda(T,a,\zeta)$ in
equation~\eqref{eq:want-field-transfer}. The activation/recovery update
\eqref{eq:want-reporter} supplies $\zeta$; the thermal step
\eqref{eq:want-thermal-step} supplies a declared temperature evolution; and
hysteresis supplies persistent $a$. Their kinetic, thermal and optical
parameters must describe the same specimen and acquisition interval.

For constant nonnegative reporter rates, its exact exponential update is a
convex combination in $[0,1]$. The scalar constant-coefficient thermal model
also has a closed exponential step. Storing these exact terms avoids repeated
Euler updates for those particular models. This improvement does not assume
rates remain constant under changing illumination or temperature. A
piecewise-constant profile must declare its interval partition, and an
unspecified nonlinear feedback network needs its own evolution definition.

The energy ledger~\eqref{eq:want-ledger} accounts for supplied lens actuation,
bridge excitation and detector amplification as well as absorbed light.
Delayed relaxation transfers energy out of an already counted store; it is
not a second optical absorption. Reporter recovery and thermochromic state
are physical memory. A logical seam involution does not reset them. A
fluorescence, resistance or thermal readout can be an auxiliary witness;
copying the primary software bit supplies no additional physical observation.

\subsection{Complex optics fits exact real-pair arithmetic}
\label{sec:want-complex}

XOP-R1 does not need a new primitive complex-number representation for the
finite calculations above. Let $S$ be a declared common real scalar sort and
represent $z=a+ib$ by $\code{Tuple}(S,S)$. Define the templates
\begin{align}
 (a,b)+(c,d)&=(a+c,b+d),\\
 (a,b)(c,d)&=(ac-bd,ad+bc),\\
 \overline{(a,b)}&=(a,-b),\qquad |(a,b)|^2=a^2+b^2,\\
 (a,b)/(c,d)&=\left(\frac{ac+bd}{c^2+d^2},
                           \frac{bc-ad}{c^2+d^2}\right),
                  \quad c^2+d^2>0,\\
 e^{i\phi}&=(\cos\phi,\sin\phi),\qquad
 e^{a+ib}=e^a(\cos b,\sin b).
 \label{eq:want-complex-pairs}
\end{align}
Each component uses existing arithmetic, sine, cosine, exponential and lazy
conditional operators with their ordinary real domains. Phase arguments are
dimensionless. A dimensioned complex amplitude is a pair of equally
dimensioned real components, not a dimension wrapper around a nonnumeric
tuple. Exact algebraic $\phi=(1+\sqrt5)/2$ must not be confused with an arbitrary
phase angle $\phi_{\mathrm{lens},\lambda}$.

A finite complex matrix $M=X+iY$ is stored as a pair of equally shaped real
matrices. Then
\begin{equation}
 M\leftrightarrow
 \begin{pmatrix}X&-Y\\Y&X\end{pmatrix},\qquad
 M^\dagger\leftrightarrow(X^{\mathsf T},-Y^{\mathsf T}).
 \label{eq:want-realification}
\end{equation}
Multiplication, conjugate transpose, trace and quadratic forms therefore
reduce to the existing finite real operators. This is an algebraic embedding;
induction on expression construction proves that it preserves the represented
complex value whenever every division and other partial operation is defined.
The frame and mode-basis annotations of each component are retained.

The real and imaginary parts of equation~\eqref{eq:want-fraunhofer} are nested
bounded real integrals of sine/cosine products. XOP's oriented finite
Riemann-integral operator can denote them when every demanded inner section
and each resulting outer integrand is Riemann-integrable on its declared
bounds. This applies separately to real and imaginary components and to each
level of a nested integral. The spectral integrals similarly
require finite wavelength bounds with zero response outside them, a proved
analytic reduction, or an explicitly declared omitted-tail approximation.
An unspecified integral over all wavelengths is not silently accepted as a
finite-integral opcode. A continuous field can be stored through a finite
function definition and parameters; arbitrary independent unknown field
values cannot be recovered from a finite expression without such structure.

Unbounded mode sums, distributional fields and a full infinite Fock-space
matrix need a separate functional representation and convergence contract.
Choosing a finite basis is legitimate when declared, but it is not an identity
between every physical field and a finite array. Algebraic expressions may
remain exact values; general transcendental terms, integrals and equality
predicates retain R11's partial-evaluation and unresolved-domain outcomes.

\subsection{The event as a quantum instrument}

For a declared finite-dimensional optical/auxiliary space, take a density
matrix $\varrho_Q\succeq0$ with $\operatorname{tr}\varrho_Q=1$ and finite
outcome operators $K_{j\ell}$ satisfying
\begin{equation}
 \sum_{j,\ell}K_{j\ell}^\dagger K_{j\ell}=I.
 \label{eq:want-kraus}
\end{equation}
The associated outcome map and probability are
\begin{equation}
 \mathcal J_j(\varrho_Q)=\sum_\ell
 K_{j\ell}\varrho_QK_{j\ell}^\dagger,
 \qquad p_j=\operatorname{tr}\mathcal J_j(\varrho_Q).
 \label{eq:want-instrument}
\end{equation}
For $p_j>0$ the conditional state is $\mathcal J_j(\varrho_Q)/p_j$; for a
zero-probability outcome it is not defined by division. The finite model must
include no-click, rejection and other possible outcomes to claim completeness.
An infinite/countable extension needs its own operator convergence conditions;
R12 does not silently serialize an infinite matrix as a finite one.

Let the recorded $j$ include the primary and secondary observations, timing
and acquisition flags. Apply the W1 classical gate
\begin{equation}
 g(j)=X\land Y\land C_t\land V\land\neg S_{\mathrm{sat}}\land\neg F,
 \qquad
 \mathcal J_e=\sum_{j:g(j)=e}\mathcal J_j.
 \label{eq:want-quantum-event}
\end{equation}
The record keeps the reason for rejection; $e=0$ is not a replacement for the
metadata that distinguishes a valid negative event from invalid acquisition.
An absent transport record is likewise not automatically an observed no-click.

\begin{proposition}[Normalized coarse-grained event instrument]
Under equation~\eqref{eq:want-kraus}, both maps $\mathcal J_e$ are completely
positive and trace non-increasing, and
$\operatorname{tr}\mathcal J_0(\varrho_Q)+
 \operatorname{tr}\mathcal J_1(\varrho_Q)=1$.
\end{proposition}
\begin{proof}
Every operator sandwich is completely positive, and finite sums preserve that
property. Cyclicity of trace and completeness give total probability one.
The two disjoint outcome groups partition this nonnegative total, so each has
trace between zero and one. The argument also holds for any positive input
with its own total trace.
\end{proof}
This is a conditional model theorem. The $K_{j\ell}$, state preparation,
dimension cutoff and detector readout are inputs to characterize. The
instrument formalism comes from measurement theory; it is not a claim that
one-bit classical output creates a qubit or quantum advantage.
\href{https://link.springer.com/article/10.1007/BF01647093}{[P3]}

NIST describes transition-edge sensors that detect photon absorption through
a resistance change near a superconducting transition. Their operating
conditions and recovery behavior belong to a specified cryogenic detector.
This supports a thermal photon witness as a real device class; it does not
establish such sensitivity for the source's unspecified ambient material.
\href{https://www.nist.gov/pml/productsservices/quantum-networks-nist/technologies-quantum-networks/single-photon-detectors}{[P4]}

\subsection{A constructive optical consequence of the pinion map}
\label{sec:want-pinion-passivity}

W1, page 18, retains the real reference map
\begin{equation}
 A_\phi=D_\phi H_4,\qquad
 D_\phi=\operatorname{diag}(\phi,\phi^{-1},1,-1),\qquad
 H_4=\frac12\begin{pmatrix}
 1&1&1&1\\1&-1&1&-1\\1&1&-1&-1\\1&-1&-1&1
 \end{pmatrix},\quad \phi=\frac{1+\sqrt5}{2}.
 \label{eq:want-pinion}
\end{equation}
This real reference map is distinct from W0's rounded and clipped integer
implementation. It is also distinct from interpreting four coordinates as
four physical field amplitudes. The following result applies only to that
last, direct power-normalized amplitude interpretation.

\begin{proposition}[Passive direct-map obstruction and scaled construction]
The unscaled $A_\phi$ cannot be the deterministic transfer of a passive
four-mode device on all power-normalized inputs. The scaled map
$B=A_\phi/\phi$ is contractive and admits an ideal real orthogonal dilation
using additional output modes.
\end{proposition}
\begin{proof}
$H_4H_4^{\mathsf T}=I$, so the singular values of $A_\phi$ are
$\phi,\phi^{-1},1,1$. In particular, the unit input
$\mathbf u=H_4^{\mathsf T}\mathbf e_1$ has
$\|A_\phi\mathbf u\|_2=\phi>1$, contradicting
equation~\eqref{eq:want-passive}. For the scaled map,
\begin{equation}
 B=\operatorname{diag}(1,\phi^{-2},\phi^{-1},-\phi^{-1})H_4.
 \label{eq:want-pinion-scaled}
\end{equation}
Every diagonal entry $d_j$ lies in $[-1,1]$. Introduce an auxiliary mode and
apply, separately for each entry,
\begin{equation}
 O_j=\begin{pmatrix}
 d_j&\sqrt{1-d_j^2}\\-\sqrt{1-d_j^2}&d_j
 \end{pmatrix},\qquad O_j^{\mathsf T}O_j=I.
 \label{eq:want-pinion-dilation}
\end{equation}
With zero classical auxiliary input amplitude, the selected output amplitude
is $d_j$ times its main input; the auxiliary output carries the remaining
power. Compose these pairs after $H_4$. Selecting the four main outputs gives
$B$, while the full enlarged transformation preserves the norm.
\end{proof}

The construction gives exact algebraic target parameters for an ideal
interferometric network. Negative amplitudes require a declared relative phase;
side-port energy must be retained in the ledger. In a quantum use, an
unexcited auxiliary port has a vacuum state and is not a deletion of its
quantum degrees of freedom. A classical zero-amplitude assignment alone is
not a complete quantum noise model.

There are now three explicit alternatives. An active amplitude device may
target $A_\phi$ with a gain and noise contract. A passive device may target
$B$, while an exact symbolic scale records that $A_\phi\mathbf u=\phi B\mathbf u$.
That identity does not recover power or remove uncertainty from measured
amplitudes; the scaled device is a distinct analog profile. Finally, a digital
photonic machine may compute the bits of coordinates and operators. Its
optical power is not the squared coordinate norm, so the direct-amplitude
obstruction does not apply to that digital encoding.

\subsection{What primary photonic-computing evidence establishes}

Bhatt, Fuller and L\'eonard reported Boolean optical logic and full adders
using coherent inputs, a spatial light modulator and intensity thresholding.
Their eight-bit addition sequence re-enters the next column's inputs through
a digital micromirror device using the previous column's result. The
demonstration uses camera/computer readout and reports power-dependent bit
errors. It establishes useful optical logic primitives, not the complete
stored, self-referential exact seed engine.
\href{https://www.nature.com/articles/s41467-026-74970-5}{[P5]}

Zhou and colleagues reported reconfigurable microring logic and two-bit
arithmetic units. Their experiment uses electrical modulation inputs, optical
outputs, electrical OR aggregation for the reported arithmetic acquisition,
and digital processing for timing and threshold recovery. This is relevant
precedent for a hybrid optical arithmetic route. Its component results are
not a system-wide latency, energy or error bound for this formalization.
\href{https://www.nature.com/articles/s41467-026-75750-x}{[P6]}

The useful contribution of W1 is therefore a substantially more explicit
\emph{photonic interface}: optical transfer, material memory, acquisition
meaning and causal feedback can be bound to exact terms. Physical execution
of every XOP operator additionally needs a device-to-bit refinement,
persistent storage, signal restoration, fan-out, clocks and a resource bound.
Those implementation obligations are addressed separately in this release's
realization contract. The optical equations above do not supply their missing
device parameters or establish a universal physical guarantee.
